\documentclass[11pt,a4paper]{article}
\usepackage[utf8]{inputenc}
\usepackage[T1]{fontenc}
\usepackage[spanish]{babel}
\usepackage{hyperref}
\usepackage{geometry}
\geometry{margin=2.5cm}
\usepackage{parskip}

\title{008 la triangulacion}
\author{Alejandro Valenzuela Rivera\\\small Universidad Autónoma de Nuevo León}
\date{\today}
\begin{document}
\maketitle

\begin{quote}
\textbf{Sesgo declarado.} Este trabajo se sostiene desde el sesgo declarado de su autor, y no finge neutralidad. Lo que afirma se puede tumbar, y ese es el punto: el falsificador está escrito, no escondido.
\end{quote}

\section*{Resumen}
Tres intuiciones con lo mundano, y ninguna se elige por parecer mejor.

\section{Por qué se sostiene}
Si solo funciona con lo complejo, la respuesta es eso es obvio.

\section{De dónde viene}
La autoridad, y el repo: la triangulación con lo mundano.

\section{El caso}
La caricatura que nadie se tomo en serio: entra por la misma puerta.

\section{Método}
Una afirmacion sostenida desde tres vertices, o con el que falta declarado.

\section{Límites, y qué lo tumba}
Una afirmacion sostenida con dos puntos y llamada triangulo.

\section{Lo que queda abierto}
\textquestiondown{}Que pasa cuando una ecuacion es a la vez el alias y no el cuerpo?

\section*{Cómo se leyó esta tesis}
Esto se escribió en el repositorio man-made-empathy, del que este candado es una pieza: candados/01 el ente y el metodo/008 la triangulacion.md. La máquina que lo compuso lo hizo desde este mismo archivo, y no agregó nada que no estuviera escrito. Si algo de esto es falso, se dice en el mismo sitio, y se corrige con fecha.

\end{document}
